\documentclass[11pt]{article}

\usepackage[margin=1in]{geometry}
\usepackage{graphicx}
\usepackage{booktabs}
\usepackage{hyperref}
\usepackage{xcolor}
\usepackage{caption}
\usepackage{float}
\usepackage{enumitem}
\usepackage{amsmath}
\usepackage{titlesec}

\hypersetup{colorlinks=true, linkcolor=blue, urlcolor=blue}
\titleformat{\section}{\large\bfseries}{\thesection}{0.6em}{}
\titleformat{\subsection}{\normalsize\bfseries}{\thesubsection}{0.6em}{}

\title{\textbf{Predicting Daily Station-Level Citi Bike Demand:\\
An End-to-End Machine Learning Pipeline}}
\author{Project 4 -- STAT GR5243 \\ Team Members: [Member 1], [Member 2], [Member 3], Ketaki Dabade}
\date{May 5, 2026}

\begin{document}
\maketitle

\begin{abstract}
\noindent We build an end-to-end machine learning pipeline to predict daily Citi Bike trip
demand at the station level in New York City. We combine Citi Bike trip records,
daily weather data, and station context features to construct one observation per
station-day. We use K-Means clustering to summarize station behavior into
interpretable groups (commuter, leisure, mixed-use, low-demand), engineer
calendar, weather, lag, and spatial features, and compare Ridge Regression,
Random Forest, Gradient Boosting, and a tuned Random Forest under a time-based
train--test split. The tuned Random Forest performs best (MAE $\approx 19.75$,
RMSE $\approx 28.13$, $R^2 \approx 0.289$ on held-out future dates), with
short-term lag demand and weather as the dominant predictors.
\end{abstract}

\section{Introduction}

Citi Bike is New York City's largest bike-share system, generating millions of
trip records each year. Reliable forecasts of station-level demand are valuable
for rebalancing operations, capacity planning, and understanding mobility
patterns across neighborhoods. This project applies the full data science
pipeline -- acquisition, cleaning, exploration, unsupervised learning, feature
engineering, supervised modeling, and evaluation -- to the problem of predicting
\emph{daily trip count per station}.

\paragraph{Predictive question.}
Given a station and a date, predict the number of trips that will originate from
that station on that day, using historical demand, weather, calendar, location,
and neighborhood context.

\paragraph{Why this is hard.}
Citi Bike demand is driven by a mix of recurring patterns (weekday commuting,
weekend leisure) and exogenous shocks (rain, cold, holidays). The distribution
of station-day demand is highly right-skewed: many low-volume stations and a
long tail of very busy stations. Models must handle nonlinearity, station
heterogeneity, and temporal dependence simultaneously.

\section{Data Acquisition and Preparation}

\subsection{Sources}
\begin{itemize}[leftmargin=1.4em]
  \item \textbf{Citi Bike trip records}: monthly trip-level CSVs from the public
        Citi Bike System Data archive (\url{https://citibikenyc.com/system-data}).
        We used 2--4 months of recent data covering both weekday and weekend
        periods.
  \item \textbf{Daily weather}: temperature, precipitation, and wind speed
        retrieved from the Open-Meteo API for the New York City coordinates
        matching the trip date range.
  \item \textbf{Station context}: station IDs, coordinates, and neighborhood
        proximity flags (e.g., \texttt{has\_park}, \texttt{has\_pier},
        \texttt{has\_plaza}, \texttt{has\_square}) constructed from station
        names and metadata.
\end{itemize}

\subsection{Cleaning}
The raw trip files were parsed and standardized:
\begin{itemize}[leftmargin=1.4em]
  \item Timestamps converted to \texttt{datetime}; trip date extracted.
  \item Trips with invalid or missing station IDs were removed.
  \item Station names standardized; coordinates verified per station.
  \item Numeric fields validated (non-negative trip durations, plausible
        coordinates).
\end{itemize}

\subsection{Aggregation}
Trips were grouped by \texttt{(station\_id, date)} to compute
\texttt{daily\_trip\_count}, then merged with daily weather (on \texttt{date})
and with station context (on \texttt{station\_id}). The final analytic table
contains one row per station per day with demand, rider/bike-type shares,
weather variables, and station context.

\section{Exploratory Data Analysis}

EDA was performed on the cleaned station-day analytic table covering
Jan 1 -- Feb 28, 2026. The figures below summarize the key patterns and the
modeling choices they motivated.

\begin{itemize}[leftmargin=1.4em]
  \item \textbf{Daily demand trend} (Figure~\ref{fig:trend}). Total system
        demand is very low ($<2{,}000$ trips/day) during the first two weeks
        of January, then jumps sharply, peaking near 48{,}000 trips on
        Jan~20 and oscillating between roughly 5{,}000 and 30{,}000 trips/day
        through February. The early-January floor reflects partial data
        coverage at the start of the window; we accounted for this by using
        a time-based train/test split and per-station lag features rather
        than naive global means.
  \item \textbf{Distribution of station-day demand}
        (Figure~\ref{fig:dist}). The histogram shows a single mode near zero
        ($\sim$4{,}000+ station-days in the lowest bin) and a long right
        tail extending past 400 trips. The boxplot confirms a median around
        40 trips/day and a heavy outlier cloud above the upper whisker
        ($\sim$175). This skew motivated the
        $\log(1+\text{daily\_trip\_count})$ target transformation.
  \item \textbf{Weekday vs.\ weekend} (Figure~\ref{fig:wkd}). Station-day
        median demand is roughly 50 on weekdays vs.\ 30 on weekends; on a
        system-wide basis, weekday averages $\approx$17{,}500 trips/day vs.\
        weekend $\approx$10{,}800 -- a $\sim$60\% weekday premium consistent
        with a commuter-dominated system. This justified including
        \texttt{is\_weekend} and cyclical day-of-week features.
  \item \textbf{Top stations} (Figure~\ref{fig:top}). The busiest stations
        cluster around midtown / Chelsea (\texttt{w 21 st \& 6 ave},
        \texttt{w 31 st \& 7 ave}, \texttt{11 ave \& w 41 st}), waterfront
        leisure (\texttt{pier 61 at chelsea piers}), and East Village /
        Union Square corridors (\texttt{lafayette st \& e 8 st},
        \texttt{broadway \& e 14 st}). The mix of commuter and leisure hubs
        motivated station clustering as a feature for the supervised model.
  \item \textbf{Weather scatters} (Figure~\ref{fig:wx}). Temperature in this
        winter window spans roughly $-15$\,$^\circ$C to $+6$\,$^\circ$C;
        demand vs.\ temperature is noisy but the highest-demand days
        ($>$30{,}000 trips) all sit at the warmer end of the range.
        Precipitation shows a clearer negative pattern: nearly all days
        with $>$10\,mm precipitation have low demand ($<$15{,}000 trips),
        while the high-demand days are concentrated at zero precipitation.
        This motivated the binary \texttt{has\_precipitation} and top-decile
        \texttt{heavy\_precipitation} flags.
  \item \textbf{Correlation heatmap} (Figure~\ref{fig:corr}).
        \texttt{daily\_trip\_count} shows a weak positive correlation with
        \texttt{member\_share} and weak negative correlations with
        \texttt{precipitation\_sum} and \texttt{max\_wind\_speed}. The
        strong $-1$ block between \texttt{casual\_share} /
        \texttt{member\_share} and between \texttt{classic\_bike\_share} /
        \texttt{electric\_bike\_share} is mechanical (the shares sum to 1)
        and does not indicate signal -- it does, however, confirm that we
        should drop one variable from each share pair if we ever use them
        as predictors.
\end{itemize}

\begin{figure}[H]
\centering
\includegraphics[width=0.7\textwidth]{Data\ Analysis/daily_demand_trend.png}
\caption{Overall daily Citi Bike demand over the study period.}
\label{fig:trend}
\end{figure}

\begin{figure}[H]
\centering
\includegraphics[width=0.48\textwidth]{Data\ Analysis/demand_distribution_histogram.png}
\includegraphics[width=0.48\textwidth]{Data\ Analysis/demand_distribution_boxplot.png}
\caption{Distribution of station-day trip counts (right-skewed).}
\label{fig:dist}
\end{figure}

\begin{figure}[H]
\centering
\includegraphics[width=0.48\textwidth]{Data\ Analysis/weekday_weekend_boxplot.png}
\includegraphics[width=0.48\textwidth]{Data\ Analysis/weekday_weekend_average_total_trips.png}
\caption{Weekday vs.\ weekend demand.}
\label{fig:wkd}
\end{figure}

\begin{figure}[H]
\centering
\includegraphics[width=0.7\textwidth]{Data\ Analysis/top_10_busiest_stations.png}
\caption{Top 10 busiest stations by total trips.}
\label{fig:top}
\end{figure}

\begin{figure}[H]
\centering
\includegraphics[width=0.48\textwidth]{Data\ Analysis/demand_vs_temperature.png}
\includegraphics[width=0.48\textwidth]{Data\ Analysis/demand_vs_precipitation.png}
\caption{Demand vs.\ temperature and precipitation.}
\label{fig:wx}
\end{figure}

\begin{figure}[H]
\centering
\includegraphics[width=0.7\textwidth]{Data\ Analysis/correlation_heatmap.png}
\caption{Correlation matrix of demand, rider/bike-type shares, and weather.}
\label{fig:corr}
\end{figure}

\section{Unsupervised Learning: Station Clustering}

To capture station heterogeneity in a compact way, we clustered \emph{stations}
(not station-days) using a profile of long-run characteristics: average and
peak demand, demand volatility (CV), casual/member share, electric-bike share,
weekend-to-weekday ratio, latitude/longitude, and neighborhood flags.

\paragraph{Method.}
Features were standardized with \texttt{StandardScaler}, and we evaluated
$k=2,\dots,8$ using the elbow method (inertia) and silhouette scores. We
selected $k=4$ as the best trade-off between interpretability and
within-cluster compactness, then visualized clusters in 2D using PCA on the
scaled feature matrix.

\paragraph{Cluster interpretation.}
The four clusters correspond roughly to:
\textbf{(0) commuter} -- high weekday demand, high member share;
\textbf{(1) low-demand} -- low average and peak demand;
\textbf{(2) leisure} -- elevated casual share and weekend ratio, often near
parks and waterfronts;
\textbf{(3) mixed-use} -- moderate demand across rider types.
The cluster label was carried into the supervised pipeline as a categorical
predictor.

\section{Feature Engineering and Preprocessing}

\subsection{Calendar features}
Month, day-of-month, day-of-week, weekend indicator, and cyclical day-of-week
encodings (\texttt{dow\_sin}, \texttt{dow\_cos}) so the model recognizes that
Sunday and Monday are adjacent in weekly time.

\subsection{Weather features}
In addition to raw temperature, precipitation, and wind, we added binary
indicators: \texttt{has\_precipitation}, \texttt{heavy\_precipitation}
(top-decile precipitation day), and \texttt{high\_wind} (top-decile wind day).
These capture nonlinear bad-weather effects that scatter plots suggested.

\subsection{Lag demand}
Per-station lag features: \texttt{lag\_1\_trip\_count} (previous-day demand)
and \texttt{lag\_7\_trip\_count} (same weekday last week). These avoid
look-ahead leakage and capture short-term persistence and weekly seasonality.
Rows missing lag values (the earliest days of each station) were dropped.

\subsection{Target transformation}
The target \texttt{daily\_trip\_count} is right-skewed; we modeled
$\log(1+y)$ and inverted predictions for evaluation. This stabilizes variance
and reduces the influence of extreme high-demand days.

\subsection{Spatial / context features}
Latitude, longitude, neighborhood flags (\texttt{has\_park},
\texttt{has\_pier}, \texttt{has\_plaza}, \texttt{has\_square}), and the
station cluster label from the unsupervised step.

\subsection{Pipeline}
A single \texttt{ColumnTransformer} applies \texttt{StandardScaler} to numeric
features and \texttt{OneHotEncoder} to the cluster label, embedded in an
\texttt{sklearn} \texttt{Pipeline} so preprocessing is fit on training data
only.

\paragraph{Excluded features.}
Rider-type and bike-type shares (\texttt{casual\_share},
\texttt{member\_share}, \texttt{classic\_bike\_share},
\texttt{electric\_bike\_share}) are derived from \emph{the same day's trips}
and would leak the target -- they were excluded from supervised inputs but
retained as inputs to the (station-level) clustering step.

\section{Supervised Modeling}

\paragraph{Task.}
Regression on $y=\log(1+\text{daily\_trip\_count})$ with the engineered
feature matrix.

\paragraph{Train/test split.}
A \emph{time-based} split holds out the most recent 20\% of dates for testing,
mirroring the real forecasting setting and preventing future-to-past leakage.

\paragraph{Models.}
\begin{enumerate}[leftmargin=1.6em]
  \item \textbf{Ridge Regression} -- linear baseline ($\alpha=1$).
  \item \textbf{Random Forest} -- 300 trees, \texttt{max\_depth=12},
        \texttt{min\_samples\_leaf=5}.
  \item \textbf{Gradient Boosting} -- 300 trees, learning rate $0.05$,
        depth 3.
  \item \textbf{Tuned Random Forest} -- selected via 3-fold
        \texttt{GridSearchCV} over \texttt{n\_estimators} $\in\{100,300\}$,
        \texttt{max\_depth} $\in\{8,12,\text{None}\}$, and
        \texttt{min\_samples\_leaf} $\in\{3,5,10\}$, scoring on negative RMSE.
        Best: \texttt{n\_estimators=300}, \texttt{max\_depth=None},
        \texttt{min\_samples\_leaf=3}.
\end{enumerate}

\section{Evaluation and Model Selection}

All metrics are reported on the \emph{original} (un-logged) trip-count scale;
predictions were exponentiated and clipped at 0.

\begin{table}[H]
\centering
\caption{Held-out test performance (sorted by RMSE).}
\label{tab:results}
\begin{tabular}{lrrr}
\toprule
\textbf{Model} & \textbf{MAE} & \textbf{RMSE} & \textbf{R\textsuperscript{2}} \\
\midrule
Tuned Random Forest & \textbf{19.75} & \textbf{28.13} & \textbf{0.289} \\
Random Forest       & 21.23          & 29.70          & 0.207 \\
Gradient Boosting   & 24.70          & 32.71          & 0.038 \\
Ridge Regression    & 25.85          & 35.78          & $-0.151$ \\
\bottomrule
\end{tabular}
\end{table}

\paragraph{Selection.}
The Tuned Random Forest dominates on every metric. We select it as the final
model for three reasons:
(i) lowest MAE/RMSE and highest $R^2$ on a time-based holdout;
(ii) it captures the nonlinear weather $\times$ calendar $\times$ location
interactions that the linear baseline cannot;
(iii) it provides interpretable feature importances for stakeholder
communication.

\paragraph{Feature importance.}
The five most important features (Table~\ref{tab:imp}) cover the three signal
families we expected: short-term demand persistence, weather, and weekly
seasonality.

\begin{table}[H]
\centering
\caption{Top features by Tuned Random Forest importance.}
\label{tab:imp}
\begin{tabular}{lr}
\toprule
\textbf{Feature} & \textbf{Importance} \\
\midrule
\texttt{lag\_1\_trip\_count}   & 0.506 \\
\texttt{mean\_temperature}     & 0.194 \\
\texttt{max\_wind\_speed}      & 0.075 \\
\texttt{precipitation\_sum}    & 0.061 \\
\texttt{lag\_7\_trip\_count}   & 0.056 \\
\texttt{dow\_sin}              & 0.020 \\
\texttt{start\_lat}            & 0.019 \\
\texttt{start\_lng}            & 0.018 \\
\texttt{heavy\_precipitation}  & 0.017 \\
\texttt{station\_cluster\_0}   & 0.011 \\
\bottomrule
\end{tabular}
\end{table}

\paragraph{Error by demand level.}
Bucketing test rows into demand quartiles
(Table~\ref{tab:bias}), the model is approximately unbiased on low to
medium-high demand but \emph{under-predicts the highest-demand days}
(actual $\approx 99.9$ vs.\ predicted $\approx 83.1$, MAE $32.6$). This is a
known consequence of the log transformation and tree-based smoothing: extreme
peaks get pulled toward the mean. For operational use, demand at the busiest
stations would benefit from a station-specific or quantile-regression
post-correction.

\begin{table}[H]
\centering
\caption{Performance by actual-demand quartile.}
\label{tab:bias}
\begin{tabular}{lrrrr}
\toprule
\textbf{Demand group} & \textbf{n} & \textbf{Avg actual} & \textbf{Avg predicted} & \textbf{MAE} \\
\midrule
Low         & 815 & 17.03 & 25.22 & 14.28 \\
Medium-Low  & 855 & 41.17 & 46.98 & 16.64 \\
Medium-High & 801 & 61.48 & 60.65 & 15.98 \\
High        & 788 & 99.90 & 83.14 & 32.62 \\
\bottomrule
\end{tabular}
\end{table}

\paragraph{Diagnostic plots.}
Predicted-vs-actual, residual, and actual-vs-predicted distribution plots
(generated in the modeling notebook) confirm the bias pattern above and show
no obvious heteroskedasticity beyond the high-demand tail.

\section{Discussion}

\paragraph{What worked.}
Lag demand alone explains roughly half of the model's predictive power, which
is consistent with the strong autocorrelation of station demand. Weather
features add meaningful lift on top of lags. The station cluster label
contributes modestly in the importance ranking but adds interpretability
without hurting accuracy.

\paragraph{Limitations.}
(i) The dataset spans only 2--4 months, so seasonality beyond a weekly cycle
is not captured. (ii) The model under-predicts peak demand; a two-stage model
(classifier for ``surge day'' + regressor) or quantile loss would help.
(iii) We did not include holiday or special-event indicators, which likely
explain some residuals.

\paragraph{Future work.}
Span a full year, add holidays/events, try LightGBM with monotonic constraints
on weather, and deploy the model behind a simple Streamlit dashboard for
station-day demand lookup.

\section{Reproducibility}

All code and data live in the project GitHub repository. The pipeline is:
\begin{enumerate}[leftmargin=1.6em]
  \item \texttt{Data Acquisition and Preprocessing/data\_acquisition\_cleaning.ipynb}
        $\to$ produces \texttt{citibike\_station\_day\_clean.csv}.
  \item \texttt{Data Analysis/5243\_project4\_analysis.py}
        $\to$ produces EDA figures and summary tables.
  \item \texttt{Project 4 Modeling.ipynb}
        $\to$ clustering, feature engineering, model training, evaluation.
\end{enumerate}
See \texttt{README.md} for environment setup and run instructions.

\section{Team Contributions}

\begin{itemize}[leftmargin=1.4em]
  \item \textbf{[Member 1]} -- Data acquisition (Citi Bike trip files,
        Open-Meteo weather, station context); data cleaning; produced the
        final aggregated \texttt{citibike\_station\_day\_clean.csv}.
  \item \textbf{[Member 2]} -- Exploratory data analysis (figures and
        summary tables in \texttt{Data Analysis/}).
  \item \textbf{[Member 3]} -- Unsupervised clustering (K-Means, PCA),
        feature engineering, supervised modeling (Ridge / Random Forest /
        Gradient Boosting / tuned RF), and in-notebook model evaluation.
  \item \textbf{Ketaki Dabade} -- Final report (this document), repository
        organization and \texttt{README}, model evaluation interpretation,
        and presentation slides.
\end{itemize}

\end{document}
